%-------------------------------------------------------------------------------
% akai-mpk-mini-play-learn
%-------------------------------------------------------------------------------
%
% \file        akai-mpk-mini-play-learn.tex
% \library     Documents
% \author      Chris Ahlstrom
% \date        2026-09-06
% \update      2026-09-07
% \version     $Revision$
% \license     $XPC_GPL_LICENSE$
%
%     This document provides LaTeX documentation for Seq66.
%
%-------------------------------------------------------------------------------

\documentclass[
 11pt,
 twoside,
 a4paper,
 final                                 % versus draft
]{article}

\input{tex/docs-structure}             % specifies document structure and layout

\usepackage{fancyhdr}
\pagestyle{fancy}
\fancyhead{}
\fancyfoot{}
\fancyheadoffset{0.005\textwidth}
\lhead{Seq66 MIDI Learn with the Akai MPK Mini Play}
\chead{}
\rhead{User Manual}
\lfoot{}
\cfoot{\thepage}
\rfoot{}
\setlength{\headheight}{14.0pt}
\makeindex

\begin{document}

\title{Seq66 v. 0.99.28 MIDI Learn with the Akai MPK Mini Play}
\author{Chris Ahlstrom \\
   (\texttt{ahlstromcj@gmail.com})}
\date{\today}
\maketitle

\clearpage                             % moves Contents to next page
\tableofcontents
\listoffigures                         % print the list of figures
\listoftables                          % print the list of tables
\parindent 0pt
\parskip 9pt
\rhead{\rightmark}         % shows section number and section name

\section{Introduction}
\label{sec:introduction}

   This document is a tutorial about using the
   MIDI Learn feature of \textsl{Seq66}.
   It supplements the "MIDI Learn" section of the user manual
   with a real case.

   The following projects support \textsl{Seq66} and its documentation (see
   the \textbf{Help} menu):

   \begin{itemize}
      \item \url{https://github.com/ahlstromcj/seq66.git}
      \item \url{https://ahlstromcj.github.io/}
      \item \url{https://ahlstromcj/github.io/docs/seq66/akai-mpk-mini-play-learn.pdf}
   \end{itemize}

\subsection{Background}
\label{subsec:background}

   The Akai MPK Mini Play Mk3 is an interesting device.
   It has a full-fledged MIDI engine, but uses an ARM processor
   that constrains what can be done with it.
   The following article explains how it works.

   \url{http://sandsoftwaresound.net/akai-mpk-mini-play-midi-implementation/}

   When attached via USB, channel 10 is used for percussion, all other
   channels respond to notes, but all program change events are
   redirection to channel 1.
   For \textsl{Seq66} and MIDI, what matters is that the MIDI
   controller portion of the emits events only on channels 1 and 10.
   Also, the ARM mini-controller filters out all SysEx other than
   its own.

\subsection{Setup}
\label{subsec:setup}

   Plug in the MPK Mini.
   Run \textsl{Seq66} and verify that it is active
   in the 
   \textbf{Edit / Preferences / MIDI Clock} and
   \textbf{MIDI Input} tabs and the
   The setup involves the following steps:

   \begin{enumber}
      \item Create a 'ctrl' file.
      \item Specify it and activate it in the 'rc' file.
      \item Set the MPK Mini as an input controller.
   \end{enumber}

\subsubsection{Setup / 'ctrl' File}
\label{subsubsec:setup_ctrl_file}

   Copy the \texttt{data/linux/qseq66.ctrl} to a file named after
   the device, such as \texttt{akai-mpk-mini-play.ctrl}. Copy this file
   to "\$home", which is usually \texttt{~/.config/seq66}.

   Read up on the \textsl{'ctrl' File} section of the user manual
   to get familiar with the control sections of this file.
   The sections relevant to MIDI Learn are

   \begin{verbatim}
      [comments]
      [loop-control]
      [mute-group-control]
      [automation-control]
   \end{verbatim}

   Only the comments section needs editing (for an explanation of
   the intended usage of the file).
   The rest are covered by the
   \textbf{Edit / Preferences / Session} tab and the
   \textbf{Learn} dialog, though the 'ctrl' can still be edited,
   if desired.

   At this point, the 'ctrl' file is in place, but not yet ready.

\subsubsection{Setup / 'rc' File}
\label{subsubsec:setup_rc_file}

   In the \textbf{Edit / Preferences / Session} tab, set up
   the new 'ctrl' file as shown here.

\begin{figure}[H]
   \centering 
   \includegraphics[scale=0.75]{midi-learning/preferences-session-select-ctrl-file.png}
   \caption*{Set New 'ctrl' File}
\end{figure}

   Note that it is marked as active and to be saved.
   Exit \textsl{Seq66}.
   In the 'rc' file \texttt{~/.config/seq66/qseq66.rc}, one should see
   the following items (adapted to one's personal setup):

\begin{verbatim}
      [comments]
      Add your comment block here

      [midi-control-file]
      active = true
      name = "akai-mpk-mini-play.ctrl"
\end{verbatim}

   Now restart \textsl{Seq66}.
   The changes made in the next section will now be able to be written
   to the 'ctrl' file.

\subsubsection{Setup / I/O Ports}
\label{subsubsec:setup_io_ports}

   Refer to the following figures.

\begin{figure}[H]
   \centering 
   \includegraphics[scale=0.75]{midi-learning/midi-clock-select-mpk.png}
   \caption*{Set MPK Mini Output Status}
\end{figure}

   Some ports might be disabled, even though present, because they
   were missing. If it is known they are plugged in click, click
   the \textbf{Re-enable} button. This also enables the input port
   entries in the \textbf{MIDI Input} tab.
   Set the MPK Mini controller as desired (usually Disabled or Off).

\begin{figure}[H]
   \centering 
   \includegraphics[scale=0.75]{midi-learning/midi-input-select-mpk.png}
   \caption*{Set MPK Mini Input Status}
\end{figure}

   Make sure the MPK Mini is enabled. Do not set \textbf{MIDI Ctrl In}
   at this time. Restart \textsl{Seq66} and verify
   that the MIDI clock and input tabs show the correct settings.

   In the 'rc' file \texttt{~/.config/seq66/qseq66.rc}, one should see
   the following items (adapted to one's personal setup):

\begin{verbatim}
      [midi-input]
       . . .
       6  1   "[6] 48:0 MPK mini Play mk3 MIDI 1"

      [midi-input-map]
       1      # map is active
       . . .
       6  1   "MPK mini Play mk3 MIDI 1"

      [midi-clock]
       . . .
       5  0   "[5] 48:0 MPK mini Play mk3 MIDI 1"

      [midi-clock-map]
       1      # map is active
       5  0   "MPK mini Play mk3 MIDI 1"
\end{verbatim}

   If not, this is an issue and should be reported.
   In the meantime, one can try to edit the 'rc' file.

\subsubsection{Setup / I/O and Automation Control}
\label{subsubsec:setup_io_and_automation_control}

   Once all the ports are configured, one can then select
   the MPK Mini as the control/automation buss.

\begin{figure}[H]
   \centering 
   \includegraphics[scale=0.75]{midi-learning/midi-input-automation-mpk.png}
   \caption*{Set MPK Mini Input as the MIDI Controller}
\end{figure}

   In the 'ctrl' file \texttt{~/.config/seq66/akai-mpk-mini-play.ctrl},
   one should see the following items (adapted to one's personal setup):

\begin{verbatim}
      [midi-control-settings]
      control-buss = "MPK mini Play mk3 MIDI 1"
      midi-enabled = true

      [midi-control-out-settings]
      output-buss = "MPK mini Play mk3 MIDI 1"
      midi-enabled = false
\end{verbatim}

\subsection{MIDI Learning}
\label{subsec:midi_learning}

TODO

\printindex

\end{document}

%-------------------------------------------------------------------------------
% vim: ts=3 sw=3 et ft=tex
%-------------------------------------------------------------------------------
